% Configuração tipográfica inspirada nas normas ABNT para trabalhos acadêmicos.
% A classe article é usada apenas como base técnica; a hierarquia do conteúdo
% segue integralmente as etapas padronizadas do relatório-base da disciplina.

\usepackage{fontspec}
\setmainfont{Times New Roman}
\setsansfont{Arial}
\setmonofont{Menlo}

\usepackage[brazil]{babel}
\usepackage{geometry}
\geometry{top=3cm,bottom=2cm,left=3cm,right=2cm}
\usepackage{microtype}
\usepackage{setspace}
\onehalfspacing

\usepackage{graphicx}
\usepackage{booktabs}
\usepackage{tabularx}
\usepackage{longtable}
\usepackage{array}
\usepackage{ragged2e}
\usepackage{caption}
\usepackage{float}
\usepackage[table]{xcolor}
\usepackage{fancyhdr}
\usepackage{etoolbox}
\usepackage[authoryear,round]{natbib}
\usepackage{hyperref}

\definecolor{cinzaescuro}{HTML}{222222}
\definecolor{cinzamedio}{HTML}{666666}
\definecolor{cinzaclaro}{HTML}{F1F1F1}
\definecolor{linhacinza}{HTML}{B7B7B7}

\hypersetup{
  colorlinks=false,
  pdfborder={0 0 0},
  pdftitle={Redução de avarias de louças sanitárias no transporte entre distribuidor e obra},
  pdfauthor={Equipe 2}
}

\setlength{\parindent}{1.25cm}
\setlength{\parskip}{0pt}
\setlength{\emergencystretch}{2em}
\clubpenalty=10000
\widowpenalty=10000
\displaywidowpenalty=10000

\makeatletter
\renewcommand\section{\@startsection{section}{1}{0pt}{1.8\baselineskip}{0.7\baselineskip}{\normalfont\normalsize\bfseries}}
\renewcommand\subsection{\@startsection{subsection}{2}{0pt}{1.4\baselineskip}{0.45\baselineskip}{\normalfont\normalsize\bfseries}}
\renewcommand\subsubsection{\@startsection{subsubsection}{3}{0pt}{1.2\baselineskip}{0.35\baselineskip}{\normalfont\normalsize\bfseries}}
\setlength{\@fptop}{0pt}
\makeatother

\captionsetup{
  font=small,
  labelfont=bf,
  justification=justified,
  singlelinecheck=false,
  labelsep=endash
}
\captionsetup[table]{position=above,skip=6pt}

\renewcommand{\arraystretch}{1.22}
\newcolumntype{Y}{>{\RaggedRight\arraybackslash}X}
\newcolumntype{L}[1]{>{\RaggedRight\arraybackslash}p{#1}}
\newcommand{\fonte}[1]{\par\vspace{2pt}\noindent{\footnotesize Fonte: #1}}
\newcommand{\notaquadro}[1]{\par\vspace{2pt}\noindent{\footnotesize Nota: #1}}
\newcommand{\pendente}[1]{\par\noindent{\small\textbf{Ponto a validar:} \textit{#1}}\par}
\newcommand{\etapafutura}[1]{\par\noindent{\small\textit{#1}}\par}

\AtBeginEnvironment{itemize}{\setlength{\itemsep}{0.2em}\setlength{\topsep}{0.4em}}
\AtBeginEnvironment{enumerate}{\setlength{\itemsep}{0.2em}\setlength{\topsep}{0.4em}}

\pagestyle{fancy}
\fancyhf{}
\fancyhead[R]{\small\thepage}
\renewcommand{\headrulewidth}{0pt}
\renewcommand{\footrulewidth}{0pt}
\setlength{\headheight}{15pt}

\newenvironment{resumorelatorio}
  {\begin{singlespace}\small\noindent}
  {\par\end{singlespace}}
